\section{legacy-risk-scan}\label{sec:risk-scan}

Repozytorium \path{letsgolegacy.legacy-risk-scan} zawiera Legacy Risk Scan: publiczną,
statyczną stronę, na którą wkleja się plik projektu albo listę zależności (\repopath{.csproj},
\repopath{packages.config}, \repopath{composer.json} i~kilka innych formatów), a~w~odpowiedzi
dostaje raport o~końcu wsparcia (end-of-life, EOL) i~znanych podatnościach frameworka oraz
każdej zależności. Raport jest liczony w~przeglądarce. Interfejs strony jest po polsku. Strona
jest publikowana przez GitHub Pages pod adresem
\url{https://konradcinkusz.github.io/letsgolegacy.legacy-risk-scan/}. Rozdział opisuje
wyłącznie technikę: co strona czyta i~jak, skąd bierze dane, jak klasyfikuje wyniki, jak
chroni plik użytkownika, jak jest testowana, budowana i~wdrażana.

\zrodla{\path{letsgolegacy.legacy-risk-scan/README.md}}

\subsection{Po co istnieje}\label{sec:risk-scan-cel}

Właściciel systemu legacy potrzebuje odpowiedzi na dwa pytania: które elementy aplikacji
straciły wsparcie producenta i~które mają znane luki bezpieczeństwa. Legacy Risk Scan
odpowiada na nie z~jednego pliku, w~kilka sekund, bez instalowania czegokolwiek i~bez
wysyłania pliku na serwer: plik jest parsowany w~przeglądarce, a~do publicznej bazy
podatności OSV trafiają wyłącznie współrzędne pakietów (ekosystem, nazwa, wersja). Dane EOL
są dołączone do strony i~datowane, więc raport mówi, jak świeże są jego dane.

Plan pracy opisuje stronę jako ticket L18 backlogu między repozytoriami. Ograniczenia
projektowe z~planu: nic z~wklejonego pliku nie opuszcza przeglądarki poza współrzędnymi
pakietów; strona nie ma analityki ani pamięci przeglądarki i~mówi o~tym; dane EOL dla
frameworków i~środowisk uruchomieniowych są dołączone i~datowane; hosting jest statyczny
(GitHub Pages), bez własnego serwera.

\zrodla{\path{letsgolegacy.legacy-risk-scan/README.md},
\path{letsgolegacy.legacy-risk-scan/docs/WORKPLAN.md}}

\subsection{Architektura i~przebieg skanu}\label{sec:risk-scan-architektura}

Strona nie ma zależności runtime: to zwykłe moduły ES serwowane w~postaci, w~jakiej są
napisane, bez bundlera, więc to, co jest przeglądane w~\path{src/}, jest dokładnie tym, co
uruchamia odwiedzający. Jedyne pakiety npm to narzędzia testowe (Playwright, axe). Silnik
(\path{src/lib/}) to czyste moduły ES, wspólne dla strony i~testów.

\begin{xltabular}{\linewidth}{@{}>{\raggedright\arraybackslash}p{4.4cm} >{\raggedright\arraybackslash}X@{}}
\toprule
Plik & Rola \\
\midrule
\endfirsthead
\toprule
Plik & Rola \\
\midrule
\endhead
\path{src/index.html}, \path{src/styles.css} & Strona i~style; polityka Content-Security-Policy w~znaczniku \repopath{meta} \\
\path{src/app.js} & Okablowanie: formularz → skan → raport; wczytanie danych EOL, przykładu i~pliku z~dysku \\
\path{src/render.js} & Raport → DOM, wyłącznie przez \repopath{textContent} i~\repopath{setAttribute} \\
\path{src/lib/manifest.js} & Tekst → wykrycie formatu → parser → manifest albo kodowany \repopath{ScanError} \\
\path{src/lib/parsers/} & \path{msbuild.js}, \path{packages-config.js}, \path{composer.js}, \path{package-json.js} \\
\path{src/lib/xml.js} & Mały, ścisły czytnik XML dla plików MSBuild i~\repopath{packages.config} \\
\path{src/lib/frameworks.js} & Target framework → środowisko, którego data końca wsparcia obowiązuje \\
\path{src/lib/versions.js} & Arytmetyka wersji: cykle wydań i~najniższa wersja dopuszczana przez ograniczenie \\
\path{src/lib/known-libraries.js} & Biblioteki JavaScript z~NuGet, pliki skryptów z~wersją w~nazwie, pakiety mapowane na produkty endoflife.date \\
\path{src/lib/eol.js} & Dopasowanie cyklu wydań i~status EOL \\
\path{src/lib/osv.js} & Klient OSV -- granica prywatności \\
\path{src/lib/report.js} & Wiersze, flagi i~liczniki podsumowania \\
\path{src/lib/servers.js} & Generacje SQL Server i~Windows Server dla formularza \\
\path{src/lib/messages.js}, \path{src/lib/format.js} & Każde zdanie widoczne dla odwiedzającego, po polsku; daty, czas względny i~liczba mnoga \\
\path{src/lib/scan.js} & Jeden skan: tekst → raport albo kodowany błąd \\
\path{scripts/} & \path{build.mjs}, \path{serve.mjs}, \path{update-eol.mjs} (z~\path{lib/eol-source.mjs}), \path{smoke-osv.mjs}, \path{wait-for-deploy.mjs}, \path{setup.sh}, skan sekretów \\
\path{data/eol.json} & Generowane dane EOL \\
\path{samples/Sklep.Legacy.csproj} & Przykład, który strona oferuje („Wypróbuj na przykładzie”) \\
\path{test/}, \path{e2e/} & Testy \repopath{node:test} z~fixture oraz Playwright \\
\bottomrule
\end{xltabular}

Przebieg jednego skanu (\repopath{runScan} w~\path{src/lib/scan.js}):

\begin{enumerate}
  \item \textbf{Parsowanie.} \repopath{parseManifest} zamienia tekst na manifest
        \texttt{\{ format, runtimes, packages, hints \}} albo rzuca \repopath{ScanError}
        z~kodem i~parametrami (sekcja~\ref{sec:risk-scan-wejscie}). Błąd wejścia kończy skan
        komunikatem; raport nie powstaje.
  \item \textbf{Współrzędne.} \repopath{coordinatesFor} wybiera każdy pakiet, który da się
        sprawdzić i~ma wersję, a~dla bibliotek JavaScript dostarczanych przez NuGet dodaje
        drugą współrzędną w~ekosystemie npm; duplikaty są usuwane.
  \item \textbf{OSV.} Jeśli są współrzędne, klient OSV pyta o~identyfikatory podatności,
        a~potem pobiera ich rekordy (sekcja~\ref{sec:risk-scan-osv}). Plik bez pakietów do
        sprawdzenia daje raport bez zapytania do OSV.
  \item \textbf{Raport.} \repopath{buildReport} łączy manifest, dane EOL, wynik OSV, datę
        dnia i~serwery wybrane w~formularzu (sekcja~\ref{sec:risk-scan-raport}).
\end{enumerate}

Degradacja zamiast porażki: gdy OSV jest nieosiągalne, raport i~tak powstaje z~danych EOL,
każdy pakiet, którego podatności nie dało się sprawdzić, jest tak oznaczony, a~strona
oferuje ponowienie. Tylko wejście, którego skan w~ogóle nie przeczyta, kończy się błędem,
a~każdy błąd -- także błąd w~kodzie -- staje się komunikatem, nigdy awarią (błąd w~kodzie jest
logowany lokalnie w~konsoli i~niczego nie wysyła). Nowy skan anuluje poprzedni
(\repopath{AbortController}); wynik anulowanego skanu nie jest raportowany. Zmiana serwerów
w~formularzu przebudowuje raport z~zapamiętanego manifestu i~wyniku OSV
(\repopath{rebuildReport}), bez parsowania i~bez sieci.

\zrodla{\path{letsgolegacy.legacy-risk-scan/README.md},
\path{letsgolegacy.legacy-risk-scan/src/lib/scan.js},
\path{letsgolegacy.legacy-risk-scan/src/app.js},
\path{letsgolegacy.legacy-risk-scan/src/lib/report.js},
\path{letsgolegacy.legacy-risk-scan/package.json}}

\subsection{Dane wejściowe: co dokładnie skanuje}\label{sec:risk-scan-wejscie}

\subsubsection{Wykrycie formatu}

Wejście to wklejony tekst albo plik wybrany z~dysku (pole pliku przyjmuje rozszerzenia
\repopath{.csproj}, \repopath{.vbproj}, \repopath{.fsproj}, \repopath{.props},
\repopath{.config}, \repopath{.json}, \repopath{.lock}). Plik z~dysku jest czytany lokalnie
(\repopath{file.text()}); plik większy niż czterokrotność limitu znaków, liczona w~bajtach,
jest odrzucany przed odczytem. \repopath{parseManifest}:

\begin{itemize}
  \item pusty tekst daje \repopath{EMPTY\_INPUT}, a~tekst dłuższy niż limit
        \path{MAX_INPUT_CHARS} (2~MB znaków, \texttt{2 * 1024 * 1024}) daje
        \path{INPUT_TOO_LARGE};
  \item usuwa BOM i~patrzy na pierwszy znak po białych znakach: \repopath{<} → XML,
        \repopath{\{} lub \repopath{[} → JSON; tekst zaczynający się od
        „Microsoft Visual Studio Solution File” → \repopath{SOLUTION\_FILE}; wszystko inne →
        \repopath{UNKNOWN\_FORMAT};
  \item XML: element główny \repopath{Project} → parser MSBuild; \repopath{packages} → parser
        \repopath{packages.config}; \repopath{configuration} → \repopath{XML\_CONFIG\_FILE}
        (wklejono \repopath{web.config} lub \repopath{app.config}); inny →
        \path{XML_UNSUPPORTED_ROOT}; błąd składni → \repopath{XML\_MALFORMED} z~przyczyną,
        wierszem i~kolumną;
  \item JSON: błąd składni → \repopath{JSON\_MALFORMED} z~wierszem i~kolumną, gdy silnik
        JavaScript podaje pozycję (obsługiwane są komunikaty V8, SpiderMonkey
        i~JavaScriptCore); nie-obiekt → \repopath{JSON\_UNSUPPORTED}; pole
        \repopath{lockfileVersion} → \path{PACKAGE_LOCK_UNSUPPORTED}
        (\repopath{package-lock.json} nie jest obsługiwany); tablica \repopath{packages} albo
        pole \repopath{content-hash} → \repopath{composer.lock}; \repopath{require} lub
        \repopath{require-dev} → \repopath{composer.json}; \repopath{dependencies},
        \repopath{devDependencies}, \path{optionalDependencies} lub \repopath{engines} →
        \repopath{package.json}; inaczej \repopath{JSON\_UNSUPPORTED}.
\end{itemize}

Każdy kod błędu ma polskie zdanie w~\path{src/lib/messages.js} (test sprawdza, że każdy kod,
który mogą zgłosić parsery, ma zdanie). Parsery zwracają tylko kody i~parametry, nigdy tekst
do wyświetlenia.

\subsubsection{Co jest czytane z~każdego formatu}

\begin{xltabular}{\linewidth}{@{}>{\raggedright\arraybackslash}p{3.4cm} >{\raggedright\arraybackslash}X@{}}
\toprule
Plik & Co skan z~niego bierze \\
\midrule
\endfirsthead
\toprule
Plik & Co skan z~niego bierze \\
\midrule
\endhead
\repopath{.csproj} / \repopath{.vbproj} w~formacie SDK & \repopath{TargetFramework} i~\repopath{TargetFrameworks}; \repopath{PackageReference} z~atrybutem \repopath{Version}, elementem \repopath{<Version>}, \repopath{VersionOverride} albo właściwością \repopath{\$(Nazwa)} zdefiniowaną w~tym samym pliku. Referencja bez wersji (central package management) jest raportowana jako „wersja nieznana” \\
\repopath{.csproj} klasyczny (sprzed .NET Core) & \path{TargetFrameworkVersion}; identyfikator i~wersja NuGet wyprowadzone z~\repopath{HintPath} do \texttt{packages\textbackslash Nazwa.Wersja\textbackslash}; biblioteki DLL spoza katalogu pakietów i~referencje GAC spoza frameworka (raportowane, nie zgadywane); pliki skryptów z~wersją w~nazwie, np. \texttt{Scripts\textbackslash jquery-1.10.2.js} \\
\path{Directory.Packages.props} & Elementy \repopath{PackageVersion} i~\path{GlobalPackageReference} \\
\repopath{packages.config} & Każdy \repopath{package} z~dokładną wersją; najnowszy \repopath{targetFramework} jako prawdopodobna wersja .NET Framework \\
\repopath{composer.json} & Ograniczenie \repopath{php} → \textbf{minimalna} dopuszczalna wersja PHP; każdy pakiet sprawdzany w~\textbf{najniższej wersji, na którą pozwala ograniczenie} (albo „nieznana”, gdy nie da się jej ustalić -- \repopath{dev-master}, \repopath{*}, \repopath{>4.0}) \\
\repopath{composer.lock} & Dokładne wersje i~ograniczenie PHP platformy \\
\repopath{package.json} & Zakresy zależności w~najniższej wersji; \repopath{engines.node} \\
\bottomrule
\end{xltabular}

\subsubsection{Pliki MSBuild}

Parser MSBuild (\path{src/lib/parsers/msbuild.js}) działa na drzewie z~\path{src/lib/xml.js};
nazwy elementów i~atrybutów porównuje bez rozróżniania wielkości liter, bo ręcznie edytowane
pliki projektu zawierają \repopath{include=} i~\repopath{<version>}.

\begin{itemize}
  \item \textbf{Format.} Projekt jest w~formacie SDK, gdy ma atrybut \repopath{Sdk}, element
        \repopath{<Sdk>} albo \texttt{<Import Sdk=...>}. Klasyczny -- gdy ma
        \path{TargetFrameworkVersion} albo przestrzeń nazw MSBuild~2003. Plik
        z~\repopath{PackageVersion} to \path{Directory.Packages.props}. Kody formatu:
        \repopath{msbuild-sdk}, \repopath{msbuild-legacy}, \repopath{cpm-props},
        \repopath{msbuild}.
  \item \textbf{Właściwości.} \repopath{\$(Nazwa)} jest rozwijane z~elementów
        \repopath{PropertyGroup} tego samego pliku, do pięciu poziomów zagnieżdżenia.
        Nierozwiązana właściwość w~wersji lub frameworku daje notatkę zamiast zgadywania.
  \item \textbf{Frameworki.} Wartości \path{TargetFrameworkVersion},
        \repopath{TargetFramework} i~\repopath{TargetFrameworks} (dzielone po średniku) są
        mapowane przez \path{src/lib/frameworks.js}: \repopath{v4.5.1}, \repopath{net451},
        \repopath{net40-client} i~\path{.NETFramework,Version=v4.8} → .NET Framework (produkt
        \repopath{dotnetfx}); \repopath{netcoreapp3.1}, \repopath{net6.0},
        \repopath{net8.0-windows} → .NET/.NET Core (\repopath{dotnet}; od wersji 5 etykieta
        „.NET”); \repopath{netstandard2.0} → .NET Standard, specyfikacja API, a~nie środowisko
        uruchomieniowe, więc bez własnej daty końca wsparcia. Wartość nierozpoznana albo
        z~nierozwiązaną właściwością dostaje przyczynę.
  \item \textbf{\repopath{PackageReference}} (\repopath{Include} albo \repopath{Update}):
        wersja z~\repopath{VersionOverride} albo \repopath{Version}, jako atrybut lub element.
        \repopath{PrivateAssets="all"} oznacza zależność deweloperską. Referencje bez wersji są
        liczone i~dają wskazówkę, żeby wkleić \path{Directory.Packages.props}.
  \item \textbf{\repopath{Reference} z~\repopath{HintPath}.} Jeśli ścieżka prowadzi przez
        katalog \repopath{packages} (także \texttt{..\textbackslash ..\textbackslash packages}
        i~\repopath{\$(SolutionDir)packages} z~nierozwiązaną właściwością), nazwa katalogu
        pakietu jest dzielona na identyfikator i~wersję
        (\path{Microsoft.AspNet.Mvc.5.2.2} → \path{Microsoft.AspNet.Mvc},
        \repopath{5.2.2}); gdy możliwych podziałów jest kilka, wygrywa ten, którego
        identyfikator równa się nazwie DLL. Wersja jest normalizowana do postaci NuGet. Ścieżka
        poza katalogiem pakietów oznacza DLL dołączoną ręcznie: wiersz „do ręcznej weryfikacji”,
        którego nie da się sprawdzić automatycznie.
  \item \textbf{\repopath{Reference} bez \repopath{HintPath}.} Zestawy frameworka
        (\repopath{System*}, \repopath{mscorlib}, \repopath{Microsoft.CSharp},
        \repopath{WindowsBase}, \repopath{PresentationCore} itd.) i~referencje bez wersji są
        pomijane i~liczone we wskazówce. Zestawy ASP.NET dostarczane poza frameworkiem
        (\repopath{System.Web.Mvc}, \repopath{System.Web.WebPages}, \repopath{System.Web.Razor},
        \repopath{System.Web.Helpers}, \repopath{System.Web.Http},
        \path{System.Web.Optimization}, \path{System.Net.Http.Formatting}) wyglądają jak
        frameworkowe, ale mają własne wersje, więc referencja GAC do nich jest raportowana.
        Inny zestaw z~wersją to komponent GAC -- wiersz do ręcznej weryfikacji.
  \item \textbf{\repopath{Content} i~\repopath{None}.} Plik skryptu z~wersją w~nazwie staje się
        pakietem npm (sekcja~\ref{sec:risk-scan-biblioteki}); wpis \repopath{packages.config}
        daje wskazówkę, że klasyczny \repopath{.csproj} wymienia tylko pakiety z~DLL, a~pełną
        listę widać w~\repopath{packages.config}. Liczba \repopath{ProjectReference} daje
        wskazówkę, żeby sprawdzić także pliki tych projektów.
  \item \textbf{Deduplikacja.} Jeden wiersz na (ekosystem, nazwa bez wielkości liter, wersja,
        możliwość sprawdzenia): kilka DLL jednego pakietu zwija się do jednego wiersza.
\end{itemize}

\subsubsection{packages.config, composer i~package.json}

\begin{itemize}
  \item \textbf{\repopath{packages.config}}: każdy \repopath{package} z~\repopath{id}
        i~dokładną \repopath{version}; \path{developmentDependency="true"} oznacza
        zależność deweloperską. Atrybuty \repopath{targetFramework} opisują framework z~chwili
        instalacji pakietu -- to dolna granica, a~nie dowód (przeniesienie projektu nie
        przepisuje pliku), więc najnowszy .NET Framework spośród nich jest prawdopodobną wersją,
        a~raport dodaje wskazówkę. Inne frameworki (\repopath{netstandard}, \repopath{net6.0})
        są wymieniane osobno.
  \item \textbf{\repopath{composer.json}}: sekcje \repopath{require} i~\repopath{require-dev}
        (ta druga -- zależności deweloperskie). Ograniczenie \repopath{php} z~\repopath{require}
        staje się wierszem środowiska: minimalną wersją PHP, na której aplikacja deklaruje
        działanie; \repopath{config.platform.php} daje dodatkowy wiersz. Wymagania platformy
        (\repopath{php-64bit} i~podobne, \repopath{hhvm}, \repopath{composer*},
        \repopath{ext-*}, \repopath{lib-*}) są pomijane i~wymienione we wskazówce. Każdy pakiet
        jest sprawdzany w~najniższej wersji dopuszczanej przez ograniczenie.
  \item \textbf{\repopath{composer.lock}}: \repopath{packages} i~\repopath{packages-dev}
        z~dokładnymi wersjami (prefiks \repopath{v} jest usuwany); gałęzie rozwojowe
        (\repopath{dev-*}, \repopath{*-dev}) mają wersję nieznaną; \repopath{platform.php}
        i~\path{platform-overrides.php} są wierszami PHP.
  \item \textbf{\repopath{package.json}}: \repopath{dependencies},
        \path{optionalDependencies} i~\repopath{devDependencies} w~najniższej wersji
        z~zakresu; \repopath{engines.node} jest wierszem środowiska Node.js. Specyfikacje spoza
        rejestru (git, ścieżki, \repopath{github:}, \repopath{user/repo}) i~etykiety
        dystrybucji (np. \repopath{latest}) mają wersję nieznaną.
\end{itemize}

\zrodla{\path{letsgolegacy.legacy-risk-scan/README.md},
\path{letsgolegacy.legacy-risk-scan/src/lib/manifest.js},
\path{letsgolegacy.legacy-risk-scan/src/lib/parsers/msbuild.js},
\path{letsgolegacy.legacy-risk-scan/src/lib/parsers/packages-config.js},
\path{letsgolegacy.legacy-risk-scan/src/lib/parsers/composer.js},
\path{letsgolegacy.legacy-risk-scan/src/lib/parsers/package-json.js},
\path{letsgolegacy.legacy-risk-scan/src/lib/frameworks.js},
\path{letsgolegacy.legacy-risk-scan/src/lib/messages.js},
\path{letsgolegacy.legacy-risk-scan/src/index.html},
\path{letsgolegacy.legacy-risk-scan/src/app.js}}

\subsection{Wersje i~ograniczenia wersji}\label{sec:risk-scan-wersje}

\path{src/lib/versions.js} odpowiada na dwa różne pytania, których nie wolno mylić: do
którego cyklu wydań należy wersja (\repopath{numericParts}, \repopath{isPrefix}) i~jaka jest
najniższa wersja, na którą pozwala ograniczenie (\repopath{lowestVersion},
\repopath{nugetLowest}). Drugie pytanie istnieje, bo \repopath{composer.json}
i~\repopath{package.json} zawierają ograniczenia, a~nie zainstalowane wersje. Gdy najniższej
wersji nie da się ustalić uczciwie, odpowiedzią jest „nieznana” z~przyczyną, nigdy
zgadywanie.

\begin{itemize}
  \item \textbf{Ograniczenia Composer i~npm} (\repopath{lowestVersion}): alternatywy
        \repopath{||} (najniższa z~nich; alternatywa gałęzi rozwojowej nie ukrywa stabilnych),
        zakres z~dywizem (\texttt{1.2 - 2.0}), operatory \repopath{>=}, \repopath{>},
        \repopath{=}, \repopath{\textasciicircum}, \repopath{\textasciitilde}; symbole
        wieloznaczne stają się zerami (\repopath{1.2.*} → \repopath{1.2.0}); \repopath{<},
        \repopath{<=} i~\repopath{!=} nie wyznaczają dolnej granicy. Wynik ma rodzaj
        \repopath{exact} (pojedyncza dokładna wersja) albo \repopath{lowest}. Przyczyny wersji
        nieznanej: \repopath{no-lower-bound}, \path{exclusive-lower-bound} (np.
        \repopath{>4.0}), \repopath{dev-branch}, \repopath{unparseable}, a~dla npm dodatkowo
        \repopath{non-registry} i~\repopath{dist-tag}.
  \item \textbf{Wersje NuGet} (\repopath{nugetLowest}): \repopath{6.0.4} -- dokładna;
        \repopath{[1.0,2.0)} -- najniższa \repopath{1.0.0}; \repopath{[1.0]} -- dokładna;
        \repopath{(1.0,)} -- nieznana (wykluczona dolna granica); \repopath{1.*} -- nieznana
        (\repopath{floating}); \repopath{\$(...)} -- nieznana
        (\repopath{unresolved-property}); brak wersji -- \repopath{not-specified}.
        Normalizacja NuGet usuwa zera wiodące, uzupełnia do trzech części, zostawia czwartą
        tylko niezerową i~usuwa metadane builda (\repopath{1.0.0.0} → \repopath{1.0.0},
        \repopath{1.01} → \repopath{1.1.0}).
  \item \textbf{Porównanie wersji}: najpierw części liczbowe, potem wersja pre-release przed
        wydaniem.
\end{itemize}

Każda przyczyna ma notatkę w~raporcie, np. „Wersja pływająca (…) -- faktycznie użyta zależy
od chwili budowania” albo „Ograniczenie … -- sprawdzono najniższą dopuszczalną wersję …”.

\zrodla{\path{letsgolegacy.legacy-risk-scan/src/lib/versions.js},
\path{letsgolegacy.legacy-risk-scan/src/lib/messages.js}}

\subsection{Biblioteki rozpoznawane po nazwie}\label{sec:risk-scan-biblioteki}

\path{src/lib/known-libraries.js} trzyma w~jednym miejscu -- „żeby dało się ją przeczytać
i~się z~nią spierać” -- całą wiedzę skanera o~konkretnych bibliotekach. Są to trzy tabele.

\textbf{Biblioteki JavaScript dostarczane przez NuGet.} Stare projekty ASP.NET dostają jQuery,
Bootstrap czy Modernizr jako pakiety NuGet, ale bazy podatności rejestrują większość ich luk
pod npm, w~rejestrze macierzystym biblioteki. Kod jest ten sam, więc takie pakiety są
sprawdzane pod obiema nazwami, a~wyniki są scalane:

\begin{tabularx}{\linewidth}{@{}>{\raggedright\arraybackslash}X >{\raggedright\arraybackslash}X@{}}
\toprule
Identyfikator NuGet (bez wielkości liter) & Pakiet npm \\
\midrule
\repopath{jquery} & \repopath{jquery} \\
\repopath{jquery.ui.combined}, \repopath{jquery.ui} & \repopath{jquery-ui} \\
\repopath{jquery.validation} & \repopath{jquery-validation} \\
\path{microsoft.jquery.unobtrusive.validation} & \path{jquery-validation-unobtrusive} \\
\path{microsoft.jquery.unobtrusive.ajax} & \path{jquery-ajax-unobtrusive} \\
\repopath{bootstrap}, \repopath{modernizr}, \repopath{lodash} & ta sama nazwa \\
\repopath{knockoutjs} & \repopath{knockout} \\
\repopath{angularjs}, \repopath{angularjs.core} & \repopath{angular} \\
\repopath{moment.js}, \repopath{momentjs} & \repopath{moment} \\
\repopath{underscore.js} & \repopath{underscore} \\
\repopath{handlebars.js} & \repopath{handlebars} \\
\bottomrule
\end{tabularx}

\textbf{Pliki skryptów z~wersją w~nazwie.} Klasyczny \repopath{.csproj} wymienia je jako
\repopath{<Content>}, co często jest jedynym śladem biblioteki. Rozpoznawane nazwy (dłuższe
i~bardziej szczegółowe najpierw): \repopath{jquery-ui}, \repopath{jquery-migrate},
\path{jquery.validate.unobtrusive} (npm \path{jquery-validation-unobtrusive}),
\repopath{jquery.validate} (npm \repopath{jquery-validation}), \repopath{jquery.signalr} (npm
\repopath{signalr}), \repopath{jquery}, \repopath{modernizr}, \repopath{knockout},
\repopath{angular}, \repopath{bootstrap}, \repopath{moment}, \repopath{lodash},
\repopath{handlebars}. Po nazwie następuje \repopath{-} lub \repopath{.} i~wersja
\repopath{x.y} albo \repopath{x.y.z}, opcjonalnie przyrostek \repopath{min}, \repopath{slim},
\repopath{slim.min}, \repopath{debug}, \repopath{intellisense} lub \repopath{vsdoc}, na końcu
\repopath{.js}. Wersja dwuczęściowa jest uzupełniana do \repopath{x.y.0}.

\textbf{Pakiety mapowane na produkty endoflife.date}, gdy wersja pakietu wyznacza cykl
wydań: npm i~NuGet \repopath{jquery} → \repopath{jquery}; npm i~NuGet \repopath{bootstrap} →
\repopath{bootstrap}; npm \repopath{angular}, NuGet \repopath{angularjs}
i~\repopath{angularjs.core} → \repopath{angularjs}; npm \repopath{@angular/core} →
\repopath{angular}; Packagist \path{laravel/framework} → \repopath{laravel}; Packagist
\path{symfony/symfony}, \path{symfony/framework-bundle}, \path{symfony/http-kernel} →
\repopath{symfony}.

\zrodla{\path{letsgolegacy.legacy-risk-scan/src/lib/known-libraries.js}}

\subsection{Dane EOL: źródło, format i~aktualizacja}\label{sec:risk-scan-eol}

\path{data/eol.json} jest \textbf{generowany, nigdy nie edytowany ręcznie}, z~publicznego API
endoflife.date (v1; dane na licencji MIT) przez \path{scripts/update-eol.mjs}. Każda
informacja o~datach wsparcia przechodzi przez funkcje \repopath{normalize*}
w~\path{scripts/lib/eol-source.mjs} -- warstwę tłumaczącą na brzegu (P11) -- więc zmiana
kształtu API psuje testy tego modułu i~job aktualizacji, a~nie dociera do odwiedzającego jako
po cichu błędna data. Plik zawiera datę pobrania (\repopath{generatedAt}); raport ją pokazuje
(„Dane o~końcu wsparcia: endoflife.date, stan na …”).

Produkty w~migawce z~2026-09-26 (wszystkie pobrane przez API v1, lista
\repopath{unavailable} pusta):

\begin{tabularx}{\linewidth}{@{}>{\raggedright\arraybackslash}p{3.2cm} >{\raggedright\arraybackslash}X >{\raggedright\arraybackslash}p{2.2cm} >{\raggedright\arraybackslash}p{2.2cm}@{}}
\toprule
Identyfikator & Etykieta & Wymagany & Cykle \\
\midrule
\repopath{dotnetfx} & .NET Framework & tak & 13 \\
\repopath{dotnet} & .NET & tak & 13 \\
\repopath{php} & PHP & tak & 18 \\
\repopath{nodejs} & Node.js & tak & 26 \\
\repopath{angular} & Angular & tak & 14 \\
\repopath{angularjs} & AngularJS & tak & 9 \\
\repopath{laravel} & Laravel & tak & 10 \\
\repopath{symfony} & Symfony & tak & 30 \\
\repopath{mssqlserver} & Microsoft SQL Server & tak & 36 \\
\repopath{windows-server} & Windows Server & tak & 20 \\
\repopath{jquery} & jQuery & nie & 4 \\
\repopath{bootstrap} & Bootstrap & nie & 4 \\
\bottomrule
\end{tabularx}

Początek pliku:

\begin{lstlisting}
{
  "schemaVersion": 1,
  "generatedAt": "2026-09-26",
  "source": {
    "name": "endoflife.date",
    "url": "https://endoflife.date",
    "api": "https://endoflife.date/api/v1/products/{product}",
    "licence": "MIT"
  },
  "products": {
    "dotnetfx": {
      "label": ".NET Framework",
      "link": "https://endoflife.date/dotnetfx",
      "api": "v1",
      "cycles": [
        {
          "name": "4.8.1",
          "label": "4.8.1",
          "releaseDate": "2022-08-09",
          "lts": false,
          "eoas": null,
          "eol": false,
          "eoes": null,
          "latest": null
        },
\end{lstlisting}

Pola cyklu: \repopath{name} i~\repopath{label} (nazwa cyklu i~etykieta, np.
\repopath{16.0} i~„2022 'Dallas'” dla SQL Server), \repopath{releaseDate},
\repopath{lts}, \repopath{eoas} (koniec aktywnego wsparcia), \repopath{eol} (koniec
wsparcia), \repopath{eoes} (koniec wsparcia rozszerzonego, ESU), \repopath{latest}.
\repopath{eol} jest datą ISO albo wartością logiczną; \repopath{eoas} i~\repopath{eoes} --
datą, wartością logiczną albo \repopath{null}, gdy produkt nie ma tej fazy. Plik ma też
\repopath{schemaVersion} (1), \repopath{source} i~listę \repopath{unavailable} produktów
opcjonalnych, których nie udało się pobrać.

\subsubsection{Generator danych}

\begin{lstlisting}
node scripts/update-eol.mjs                          fetch, validate, write
node scripts/update-eol.mjs --keep-date-if-unchanged keep the old date when nothing
                                                     but the date would change
node scripts/update-eol.mjs --out <file> --date YYYY-MM-DD
\end{lstlisting}

\begin{itemize}
  \item Domyślnie \repopath{-{}-out} to \path{data/eol.json}, a~\repopath{-{}-date} -- dzisiejsza
        data UTC. Ręcznie: \texttt{npm run update-eol}; za proxy HTTP z~prefiksem
        \path{NODE_USE_ENV_PROXY=1} (Node~22.21 lub nowszy).
  \item Każdy produkt jest pobierany z~\path{https://endoflife.date/api/v1/products/<id>}
        z~nagłówkami \texttt{Accept: application/json} i~\repopath{User-Agent}
        identyfikującym odświeżanie danych, z~limitem 30~sekund. Odpowiedzi 429, 5xx i~błędy
        sieci są ponawiane (do trzech prób, przerwy 2~s i~4~s); 404 oznacza, że produkt nie
        istnieje.
  \item Odpowiedź v1 o~nieoczekiwanym kształcie daje ostrzeżenie i~przejście na starsze API v0
        (\path{https://endoflife.date/api/<id>.json}); plik zapisuje, które API użyto dla
        każdego produktu. Normalizacja v0 odwraca flagi wsparcia („\texttt{support: true}”
        znaczy, że aktywne wsparcie trwa, czyli \repopath{eoas} = \repopath{false}).
  \item Plik jest zapisywany tylko wtedy, gdy pobrano każdy wymagany produkt i~wynik przeszedł
        walidację -- częściowy albo zniekształcony zbiór nigdy nie zastępuje kompletnego. Brak
        wymaganego produktu albo błąd walidacji kończy skrypt kodem 1 bez zmiany pliku;
        niedostępny produkt opcjonalny trafia na listę \repopath{unavailable}.
  \item Walidacja (\repopath{validateDataset}): każdy wymagany produkt jest obecny; każdy cykl
        ma nazwę, \repopath{eol} jest wartością logiczną lub datą ISO, \repopath{eoas}
        i~\repopath{eoes} -- \repopath{null}, wartością logiczną lub datą; sondy, czyli
        wersje, które na pewno istnieją, muszą trafić w~cykl: .NET Framework 4.5.1, 4.6.2
        i~4.8; .NET 3.1, 6.0 i~8.0; PHP 7.4.0 i~8.2.0; Node.js 14.0.0 i~20.0.0; Angular
        12.0.0; Laravel 8.0.0; Symfony 4.4.0. SQL Server musi mieć cykl 2016, Windows Server
        -- 2012, a~\repopath{generatedAt} musi być datą ISO.
  \item \path{--keep-date-if-unchanged} zostawia starą datę, gdy poza datą nic się nie
        zmieniło (porównanie treści pomija \repopath{generatedAt}).
  \item Skrypt drukuje tabelę wybranych dat (m.in. .NET Framework 4.5.2--4.8, .NET 6, 8, 9,
        10, SQL Server 2014--2019, Windows Server 2012--2019, PHP 7.4, 8.1, 8.2) do przejrzenia
        w~pull requeście -- drukowaną, nigdy asertowaną, bo źródłem rozstrzygającym jest
        endoflife.date -- także do podsumowania kroku w~GitHub Actions.
\end{itemize}

Odświeżanie danych automatyzuje workflow \path{update-eol.yml}
(sekcja~\ref{sec:risk-scan-ci}): co miesiąc regeneruje plik, uruchamia na nim testy
jednostkowe i~otwiera pull request; pull request zmieniający generator dostaje dane
zregenerowane na własnej gałęzi, więc generator i~dane nie mogą się rozjechać.

\zrodla{\path{letsgolegacy.legacy-risk-scan/README.md},
\path{letsgolegacy.legacy-risk-scan/data/eol.json},
\path{letsgolegacy.legacy-risk-scan/scripts/update-eol.mjs},
\path{letsgolegacy.legacy-risk-scan/scripts/lib/eol-source.mjs}}

\subsection{Reguły oceny: cykl wydań i~status EOL}\label{sec:risk-scan-reguly}

\subsubsection{Dopasowanie cyklu wydań}

Wersja jest dopasowywana do cyklu wydań po najdłuższym prefiksie liczbowym
(\repopath{findCycle} w~\path{src/lib/eol.js}): \repopath{4.5.1} → cykl \repopath{4.5.1},
\repopath{8.0} → \repopath{8}, \repopath{7.1.3} → \repopath{7.1}; \repopath{4.8.1} nie jest
raportowane jako \repopath{4.8}. Krótka nazwa cyklu, która ma bardziej szczegółowe
rodzeństwo, oznacza swoje wydanie \repopath{.0}: na liście .NET Framework \repopath{4} to 4.0
(4.5 i~4.8 są osobnymi cyklami), więc niewymieniona wersja 4.7.2 nie zostanie zgłoszona jako
4.0; bez rodzeństwa \repopath{8} oznacza całe 8.x.

endoflife.date nie wymienia każdego wydania (np. Laravel 5.6 leży między wymienionymi 5.5
i~5.8, a~stare wydania z~czasem znikają). Dla wersji niewymienionej zwracany jest najbliższy
\emph{nowszy} wymieniony cykl jako punkt odniesienia: jeśli nawet on stracił wsparcie, ta
starsza wersja też je straciła (status \repopath{eol} z~notatką „Tej wersji nie ma
w~endoflife.date, ale nowsza wersja … straciła wsparcie …”); jeśli nie stracił -- odpowiedzią
jest „brak danych”, a~nie zgadywanie. W~szczególności wersja starsza od wszystkich
wymienionych cykli jest po końcu wsparcia, jeśli najstarszy wymieniony cykl już go nie ma.

\subsubsection{Status cyklu}

Status jest liczony z~dat w~dniu tworzenia raportu, a~nie zamrażany przy pobraniu danych:
według komentarza w~kodzie koniec wsparcia .NET~8 przypadający na 2026-11-10 zmienia wiersz
z~„koniec wsparcia w~ciągu 12 miesięcy” na „koniec wsparcia” tego dnia, nawet jeśli nikt nie
odświeżył danych. Data danych jest pokazywana w~raporcie, więc ich świeżość nigdy nie jest
domniemana.

\begin{tabularx}{\linewidth}{@{}>{\raggedright\arraybackslash}p{4.8cm} >{\raggedright\arraybackslash}p{2.4cm} >{\raggedright\arraybackslash}X@{}}
\toprule
Warunek (\repopath{eol} cyklu) & Status & Etykieta w~raporcie \\
\midrule
data ≤ dziś (sam dzień końca się liczy) & \repopath{eol} & „Koniec wsparcia” \\
data ≤ dziś + 12 miesięcy & \repopath{eol-soon} & „Koniec wsparcia w~ciągu roku” \\
data późniejsza & \repopath{supported} & -- \\
\repopath{true} (zakończone, bez daty) & \repopath{eol} & „Koniec wsparcia” \\
\repopath{false} (nie ogłoszono końca) & \repopath{supported} & notatka „Producent nie ogłosił daty końca wsparcia.” \\
wersja bez dopasowanego cyklu & \repopath{no-data} & notatka „Brak danych o~cyklu życia tej wersji w~endoflife.date.” \\
\bottomrule
\end{tabularx}

Dodawanie miesięcy przycina dzień do ostatniego dnia miesiąca (2028-02-29 + 12 → 2029-02-28).
Dodatkowe notatki: dla cyklu po końcu wsparcia z~datą \repopath{eoes} -- czy płatne
rozszerzone aktualizacje bezpieczeństwa (ESU) są jeszcze dostępne i~do kiedy, czy już się
skończyły; dla cyklu wspieranego, którego \repopath{eoas} minęło -- że od tej daty wydawane są
już tylko poprawki bezpieczeństwa.

\subsubsection{Środowiska i~serwery}

\begin{itemize}
  \item .NET Standard jest oceniony (\repopath{assessed}) bez daty EOL, z~notatką, że to
        specyfikacja API, a~nie środowisko uruchomieniowe.
  \item Środowisko wyznaczone ograniczeniem (\repopath{composer.json},
        \repopath{engines.node}) jest pokazywane jako „co najmniej …” z~notatką, że na serwerze
        może działać nowsza wersja; wersja PHP z~\repopath{config.platform} ma notatkę
        o~konfiguracji platformy Composera.
  \item SQL Server i~Windows Server nie pochodzą z~pliku: odwiedzający może wybrać je
        w~formularzu. endoflife.date wymienia SQL Server per service pack
        (\repopath{13.0-sp3}, etykieta „2016 SP3”), a~Windows Server per kanał, łącznie
        z~krótko żyjącymi kanałami półrocznymi, więc formularz (\repopath{serverOptions}) oferuje
        jedną pozycję na generację produktu, wskazującą cykl wspierany najdłużej -- ostatni
        service pack. Pomijane są kanały \repopath{-sac}, \repopath{-ac}, \repopath{-acp}
        i~pakiety „Azure Connect”; etykiety tracą nazwy kodowe, prefiks produktu i~oznaczenia
        kanałów; generacje są sortowane od najnowszego pierwszego wydania. Wiersz serwera
        w~raporcie ma notatkę, że terminy dotyczą ostatniego service packa, a~instalacje ze
        starszym straciły wsparcie wcześniej.
\end{itemize}

\zrodla{\path{letsgolegacy.legacy-risk-scan/README.md},
\path{letsgolegacy.legacy-risk-scan/src/lib/eol.js},
\path{letsgolegacy.legacy-risk-scan/src/lib/servers.js},
\path{letsgolegacy.legacy-risk-scan/src/lib/report.js},
\path{letsgolegacy.legacy-risk-scan/src/lib/messages.js}}

\subsection{Podatności: klient OSV}\label{sec:risk-scan-osv}

\path{src/lib/osv.js} jest jedynym kodem strony, który łączy się z~czymkolwiek poza własnym
originem strony, a~funkcja \repopath{toQueries()} jest jedynym miejscem, które decyduje, co
zostaje wysłane: współrzędne pakietów -- ekosystem, nazwa, wersja -- i~nic więcej. Rekordy OSV
są tu redukowane do kilku pól, które pokazuje raport (tłumaczenie na brzegu, P11), więc zmiana
schematu OSV psuje ten plik i~jego testy, a~nie interfejs.

\begin{itemize}
  \item \textbf{Zapytanie.} \texttt{POST https://api.osv.dev/v1/querybatch} z~ciałem
        \texttt{\{"queries": [...]\}}, w~którym każde zapytanie ma postać
        \texttt{\{ package: \{ ecosystem, name \}, version \}}.
        Współrzędne są deduplikowane; paczka ma najwyżej 1000 zapytań (udokumentowany limit
        OSV), a~wyniki z~\repopath{next\_page\_token} są dociągane kolejnymi zapytaniami, do
        dziesięciu stron. Nazwy są wysyłane dokładnie tak, jak są zapisane: OSV dopasowuje
        identyfikatory NuGet z~rozróżnianiem wielkości liter.
  \item \textbf{Rekordy.} Dla każdego zwróconego identyfikatora jedno \texttt{GET
        /v1/vulns/<id>}, najwyżej cztery równolegle. Rekord, którego nie udało się pobrać, jest
        \repopath{null}: raport pokazuje wtedy identyfikator i~link bez opisu, zamiast przerywać
        skan.
  \item \textbf{Żądania} mają limit 20~sekund, \texttt{credentials: 'omit'},
        \texttt{referrerPolicy: 'no-referrer'} i~\texttt{cache: 'no-store'}. Błędy są
        klasyfikowane jako \repopath{network}, \repopath{timeout}, \repopath{http} (ze statusem)
        i~\repopath{contract} (odpowiedź nie jest JSON albo liczba wyników nie zgadza się
        z~liczbą zapytań). Anulowany skan kończy się przyczyną anulowania, a~nie błędem OSV.
  \item \textbf{Streszczenie podatności.} Waga z~\path{database_specific.severity} rekordu
        albo wpisu \repopath{affected} (LOW → niska, MODERATE i~MEDIUM → średnia, HIGH →
        wysoka, CRITICAL → krytyczna); aliasy; \repopath{summary} albo pierwszy akapit
        \repopath{details} (skracany do 200 znaków); wersje z~poprawką (\repopath{fixed}) tylko
        z~wpisów \repopath{affected} dla tego pakietu i~tylko wyższe od używanej; link
        \path{https://osv.dev/vulnerability/<id>}.
  \item \textbf{Scalanie.} Podatności znalezione pod kilkoma nazwami (NuGet \repopath{jQuery}
        i~npm \repopath{jquery}) są scalane: wycofane (\repopath{withdrawn}) odpadają,
        duplikaty są zwijane po identyfikatorze i~wspólnym aliasie, kolejność -- od najwyższej
        wagi.
\end{itemize}

\zrodla{\path{letsgolegacy.legacy-risk-scan/src/lib/osv.js},
\path{letsgolegacy.legacy-risk-scan/README.md}}

\subsection{Raport: wiersze, statusy i~podsumowanie}\label{sec:risk-scan-raport}

\path{src/lib/report.js} jest czysty: bez DOM, bez sieci i~bez tekstu do wyświetlenia. Każdy
wiersz niesie kody notatek, które \path{src/lib/messages.js} zamienia na polskie zdania, więc
klasyfikację da się testować osobno.

Wiersz ma rodzaj \repopath{runtime} (framework lub środowisko z~pliku),
\repopath{infrastructure} (serwer z~formularza) albo \repopath{package}, oraz pola
\repopath{name}, \repopath{version}, \repopath{spec}, \repopath{versionKind},
\repopath{source} (skąd pochodzi wiersz, np. \repopath{HintPath}), \repopath{ecosystem},
\repopath{eol}, \repopath{advisoryState}, \repopath{advisories}, \repopath{assessed},
\repopath{notes}, \repopath{flags} i~\repopath{status}. Flagi: \repopath{eol} -- wsparcie
cyklu się skończyło; \repopath{eolSoon} -- kończy się w~ciągu 12 miesięcy;
\repopath{vulnerable} -- co najmniej jedna znana, niewycofana podatność. Wiersz jest
\repopath{assessed}, gdy wiadomo dość, by powiedzieć „nic nie znaleziono”: pakiet ma wersję
i~odpowiedź OSV, środowisko -- dopasowany cykl.

\begin{tabularx}{\linewidth}{@{}>{\raggedright\arraybackslash}p{2.4cm} >{\raggedright\arraybackslash}p{5.2cm} >{\raggedright\arraybackslash}X@{}}
\toprule
Status & Etykieta & Kiedy \\
\midrule
\repopath{eol} & „Koniec wsparcia” & flaga \repopath{eol} \\
\repopath{vulnerable} & „Znane podatności” & flaga \repopath{vulnerable}, bez \repopath{eol} \\
\repopath{eol-soon} & „Koniec wsparcia w~ciągu roku” & flaga \repopath{eolSoon}, bez poprzednich \\
\repopath{unknown} & „Nieustalone” & brak flag, wiersz nieoceniony \\
\repopath{ok} & „Bez zastrzeżeń” & brak flag, wiersz oceniony \\
\bottomrule
\end{tabularx}

Status główny to najgorsza flaga, w~kolejności z~tabeli. Wiersz może liczyć się w~kilku
kategoriach podsumowania (środowisko po końcu wsparcia ze znanymi podatnościami jest w~obu),
co strona mówi wprost. Stan sprawdzenia podatności pakietu (\repopath{advisoryState}):
\repopath{checked}, \repopath{not-checkable} (DLL dołączona ręcznie, komponent GAC --
„nie da się sprawdzić automatycznie”), \repopath{unknown-version} („nie sprawdzono --
nieznana wersja”), \repopath{error} („nie sprawdzono -- brak połączenia z~OSV”),
\repopath{not-checked}; dla środowisk i~serwerów \repopath{not-applicable}.

Wynik \repopath{buildReport} zawiera format pliku, datę raportu, datę danych EOL, stan OSV
(i~rodzaj błędu), wskazówki do pliku, dwie listy -- \repopath{platform} (frameworki,
środowiska, serwery) i~\repopath{dependencies} (pakiety), każda sortowana od najgorszego
statusu, a~potem po nazwie z~polskim porządkiem -- oraz podsumowanie: liczba pozycji, po końcu
wsparcia, z~końcem w~ciągu 12 miesięcy, ze znanymi podatnościami, bez zastrzeżeń,
nieustalonych i~liczba różnych podatności.

Strona renderuje raport (\path{src/render.js}) wyłącznie przez \repopath{textContent}
i~\repopath{setAttribute}, nigdy \repopath{innerHTML}: streszczenia podatności są tekstem
z~zewnątrz, a~wklejony plik należy do odwiedzającego, więc żadne z~nich nie może zostać
zinterpretowane jako znaczniki. Wiersz pokazuje do trzech podatności, resztę za przyciskiem
„Pokaż pozostałe (…)”. Raport ma krótkie objaśnienie każdego poziomu ryzyka, sekcje
„Platforma”, „Zależności” i~„Uwagi do pliku”, informację o~źródłach i~datach danych oraz
przycisk „Drukuj lub zapisz jako PDF” (\repopath{window.print()}). Gdy OSV jest nieosiągalne,
raport pokazuje ostrzeżenie i~przycisk „Sprawdź podatności ponownie”.

\zrodla{\path{letsgolegacy.legacy-risk-scan/src/lib/report.js},
\path{letsgolegacy.legacy-risk-scan/src/lib/messages.js},
\path{letsgolegacy.legacy-risk-scan/src/render.js},
\path{letsgolegacy.legacy-risk-scan/src/index.html}}

\subsection{Prywatność i~bezpieczeństwo strony}\label{sec:risk-scan-prywatnosc}

Prywatność jest właściwością konstrukcji, a~nie deklaracją:

\begin{itemize}
  \item Wklejony plik jest parsowany \textbf{w~przeglądarce}; nic nie jest wysyłane.
  \item Połączenia wychodzące idą wyłącznie do publicznego API OSV.dev: zapytanie
        \repopath{querybatch}, którego ciało buduje jedna funkcja (\repopath{toQueries()}),
        a~potem po jednym \repopath{GET} na każdą podatność zwróconą przez OSV, po
        identyfikatorze, który podało OSV. Poza tym strona pobiera tylko własne pliki
        (\path{data/eol.json}, przykład).
  \item Polityka Content-Security-Policy strony pozwala łączyć się tylko ze stroną
        i~z~\repopath{api.osv.dev} i~nie pozwala ładować skryptów, stylów ani fontów skądinąd; strona
        nie ma skryptów ani stylów inline, więc nie da się ich też wstrzyknąć. Ustawiona jest
        też \repopath{referrer} \repopath{no-referrer}.
  \item Brak analityki, ciasteczek i~pamięci przeglądarki.
\end{itemize}

\begin{lstlisting}
<meta http-equiv="Content-Security-Policy" content="default-src 'none'; script-src 'self'; style-src 'self'; img-src 'self' data:; connect-src 'self' https://api.osv.dev; base-uri 'none'; form-action 'none'; manifest-src 'none'">
<meta name="referrer" content="no-referrer">
\end{lstlisting}

Test e2e prywatności sprawdza to na zbudowanej stronie: że wszystkie żądania zewnętrzne idą
wyłącznie do \path{https://api.osv.dev}; że ciało każdego \repopath{POST} ma tylko klucz
\repopath{queries}, a~każde zapytanie tylko \repopath{package} (z~\repopath{ecosystem}
i~\repopath{name}) i~\repopath{version}; że ciało nie zawiera nazwy projektu, nazwy ręcznie
dołączonej biblioteki DLL ani ścieżek (znaku \repopath{\textbackslash}); że żądania nie mają
nagłówka \repopath{Referer}; że \repopath{GET} idą tylko do \path{/v1/vulns/<id>}; że polityka
CSP ma oczekiwane \repopath{connect-src} i~nie zgłosiła naruszeń (licząc także próby
zablokowane, zanim powstało żądanie); że ciasteczka, \repopath{localStorage},
\repopath{sessionStorage} i~IndexedDB są puste. Sekcja strony „Szczegóły techniczne” mówi też
wprost, że GitHub Pages i~\repopath{api.osv.dev} -- jak każdy serwer -- widzą adres IP
połączenia.

Pozostałe zabezpieczenia: czytnik XML (\path{src/lib/xml.js}) nie obsługuje deklaracji encji
-- dekoduje tylko pięć encji predefiniowanych i~odwołania numeryczne, więc nie ma rozwijania
encji (bez XXE i~„billion laughs”), a~DOCTYPE jest pomijany bez interpretacji; napisano go
ręcznie, bo parsery muszą być czystymi funkcjami działającymi identycznie w~przeglądarce
i~pod \texttt{node -{}-test}, a~strona nie dostarcza kodu stron trzecich. Szablon pull requestu
wymaga zaznaczenia, czy zmiana dotyka tego, co strona wysyła przez sieć, i~jeśli tak --
potwierdzenia, że test prywatności nadal asertuje ciało żądania. Szablon zgłoszenia błędu
prosi, żeby nie wklejać prawdziwego pliku projektu klienta, tylko zredukowany fragment bez
nazw, ścieżek i~wewnętrznych feedów.

\zrodla{\path{letsgolegacy.legacy-risk-scan/README.md},
\path{letsgolegacy.legacy-risk-scan/src/index.html},
\path{letsgolegacy.legacy-risk-scan/src/lib/osv.js},
\path{letsgolegacy.legacy-risk-scan/src/lib/xml.js},
\path{letsgolegacy.legacy-risk-scan/e2e/scan.spec.mjs},
\path{letsgolegacy.legacy-risk-scan/.github/pull_request_template.md},
\path{letsgolegacy.legacy-risk-scan/.github/ISSUE_TEMPLATE/bug.yml}}

\subsection{Budowanie i~uruchamianie lokalne}\label{sec:risk-scan-lokalnie}

Wymagany jest Node~22 lub nowszy (\repopath{engines.node} w~\path{package.json}: \texttt{>=22};
test runner i~\repopath{fetch}). Polecenia z~README:

\begin{lstlisting}[escapechar=!]
./scripts/setup.sh           # checks Node, runs npm ci, installs the secret-scanning hook
npm start                    # build dist/ and serve it at
                             # http://127.0.0.1:4173/letsgolegacy.legacy-risk-scan/
npm test                     # unit tests (node:test)
npx playwright install chromium   # once, for the browser tests
npm run e2e                  # build, then Playwright against dist/ with OSV mocked
npm run smoke:osv            # the real OSV API: CORS for the published origin + contract
npm run e2e:live             # the built page against the real OSV API
BASE_URL=https://konradcinkusz.github.io/letsgolegacy.legacy-risk-scan/ npm run e2e:live
                             # !\ldots!or against the published page
./scripts/scan-secrets.sh    # the same gitleaks scan CI runs
\end{lstlisting}

\begin{xltabular}{\linewidth}{@{}>{\raggedright\arraybackslash}p{3.7cm} >{\raggedright\arraybackslash}X@{}}
\toprule
Skrypt & Działanie \\
\midrule
\endfirsthead
\toprule
Skrypt & Działanie \\
\midrule
\endhead
\path{scripts/build.mjs} (\texttt{npm run build}) & Składa statyczną stronę w~\path{dist/}: zawartość \path{src/}, \path{data/eol.json} i~\path{samples/Sklep.Legacy.csproj}, bez bundlowania i~transpilacji. Zapisuje \path{dist/build.json} z~commitem (\repopath{GITHUB\_SHA} albo \texttt{git rev-parse HEAD}, inaczej \repopath{unknown}) i~czasem budowy; kontrola po wdrożeniu czeka, aż opublikowana strona poda ten commit \\
\path{scripts/serve.mjs} & Mały serwer statyczny dla \path{dist/} (używany przez \texttt{npm start} i~testy e2e): \texttt{node scripts/serve.mjs [-{}-dir dist] [-{}-port 4173] [-{}-base /letsgolegacy.legacy-risk-scan/]}. Domyślnie \path{dist}, port ze zmiennej \repopath{PORT} albo 4173, baza \repopath{/}; nasłuch na \repopath{127.0.0.1}. \repopath{-{}-base} montuje stronę pod ścieżką, tak jak GitHub Pages serwuje stronę projektu, więc bezwzględny URL, który zepsułby się po publikacji, psuje się najpierw lokalnie. Nagłówki \texttt{Cache-Control: no-store} i~\texttt{X-Content-Type-Options: nosniff}; adres bazy bez ukośnika → 301; poza bazą → 404; próba wyjścia poza katalog → 403; katalog → \path{index.html} \\
\path{scripts/setup.sh} & Sprawdza Node (co najmniej 22), uruchamia \texttt{npm ci} (tylko narzędzia deweloperskie), instaluje hook pre-commit skanujący sekrety (bez gitleaks i~Dockera hook odmawia commitów -- celowo, P5) i~podaje, które opcjonalne narzędzie włącza którą warstwę testów \\
\path{scripts/update-eol.mjs} (\texttt{npm run update-eol}) & Odświeżenie \path{data/eol.json} (sekcja~\ref{sec:risk-scan-eol}) \\
\path{scripts/smoke-osv.mjs} (\texttt{npm run smoke:osv}) & Test dymny prawdziwego API OSV (sekcja~\ref{sec:risk-scan-testy}) \\
\path{scripts/wait-for-deploy.mjs} & Czekanie, aż opublikowana strona poda dany commit (sekcja~\ref{sec:risk-scan-deploy}) \\
\path{scripts/scan-secrets.sh} & Lokalne odbicie skanu sekretów z~CI: domyślnie pełna historia, \repopath{-{}-staged} (to, co skanuje hook), \repopath{-{}-dir} (drzewo robocze jako pliki) \\
\bottomrule
\end{xltabular}

\zrodla{\path{letsgolegacy.legacy-risk-scan/README.md},
\path{letsgolegacy.legacy-risk-scan/package.json},
\path{letsgolegacy.legacy-risk-scan/scripts/build.mjs},
\path{letsgolegacy.legacy-risk-scan/scripts/serve.mjs},
\path{letsgolegacy.legacy-risk-scan/scripts/setup.sh},
\path{letsgolegacy.legacy-risk-scan/scripts/scan-secrets.sh}}

\subsection{Testy}\label{sec:risk-scan-testy}

\begin{xltabular}{\linewidth}{@{}>{\raggedright\arraybackslash}p{3.3cm} >{\raggedright\arraybackslash}X >{\raggedright\arraybackslash}p{3.5cm}@{}}
\toprule
Warstwa & Co & Kiedy \\
\midrule
\endfirsthead
\toprule
Warstwa & Co & Kiedy \\
\midrule
\endhead
Jednostkowe (\path{test/}) & Parsery, ograniczenia wersji, dopasowanie EOL, klient OSV, raport, polskie teksty; zatwierdzony \path{data/eol.json} & Każdy PR i~push; w~miesięcznym odświeżeniu danych przed otwarciem PR \\
E2E z~makietą OSV (\path{e2e/scan.spec.mjs}) & Przykład daje raport z~wierszem EOL i~podatnością; prywatność; uszkodzone wejście; OSV niedostępne i~ponowienie; wybór serwerów; plik z~dysku; klawiatura; szerokość 390~px; axe w~trybie jasnym i~ciemnym & Każdy PR i~push; przed każdym wdrożeniem \\
Dymne, prawdziwe OSV (\path{scripts/smoke-osv.mjs}, \path{e2e/live.spec.mjs}) & Preflight CORS i~POST dla opublikowanego originu; podatność przykładu; strona w~prawdziwej przeglądarce & Każdy PR i~push, co tydzień; błąd sieci jest ostrzeżeniem, zła odpowiedź -- porażką \\
Po wdrożeniu (\path{deploy.yml} → \repopath{verify}) & Opublikowana strona, prawdziwe OSV, zrzut ekranu jako artefakt & Po każdym wdrożeniu \\
\bottomrule
\end{xltabular}

\subsubsection{Testy jednostkowe}

\texttt{npm test} uruchamia \texttt{node -{}-test -{}-test-reporter=spec "test/**/*.test.mjs"}.
Pliki testów: \path{xml.test.mjs} (elementy, atrybuty, CDATA, komentarze, deklaracja
i~DOCTYPE; brak rozwijania deklarowanych encji), \path{manifest.test.mjs} (wykrycie formatu,
odrzucenie zbyt dużego wejścia, obcięty plik XML i~JSON z~pozycją błędu, pozycje błędów JSON
różnych silników; „nigdy nie pada”: każde obcięcie każdego obsługiwanego pliku albo się
parsuje, albo kończy się \repopath{ScanError}), \path{msbuild.test.mjs} (klasyczny
\repopath{.csproj} przykładu, \repopath{HintPath}, zwijanie DLL jednego pakietu, DLL ręczne
i~GAC, pliki skryptów, zestawy frameworka i~spoza frameworka, projekt SDK, central package
management, \path{Directory.Packages.props}), \path{other-parsers.test.mjs}
(\repopath{packages.config}, \repopath{composer.json}, \repopath{composer.lock},
\repopath{package.json}), \path{frameworks.test.mjs}, \path{versions.test.mjs},
\path{eol.test.mjs} (reguły dopasowania cyklu i~statusu, ESU, arytmetyka kalendarza),
\path{eol-source.test.mjs} (normalizacja v1 i~v0, walidacja zbioru, porównanie treści),
\path{eol-data.test.mjs}, \path{osv.test.mjs} (dokładne ciało zapytania, deduplikacja, brak
poświadczeń i~referrera, podział na paczki, stronicowanie, klasyfikacja błędów, anulowanie,
brakujące rekordy, streszczenia, scalanie), \path{report.test.mjs}, \path{scan.test.mjs},
\path{servers.test.mjs}, \path{messages.test.mjs} (każdy kod błędu, przyczyny wersji
i~notatki ma polskie zdanie; liczba mnoga; daty). Fixture: \path{test/fixtures/manifests/},
nagrane rekordy OSV w~\path{test/fixtures/osv/} (\repopath{GHSA-5crp-9r3c-p9vr},
\repopath{GHSA-gxr4-xjj5-5px2}, \repopath{GHSA-rmxg-73gg-4p98}) i~odpowiedzi endoflife.date
w~\path{test/fixtures/endoflife/}.

\path{eol-data.test.mjs} sprawdza zatwierdzony \path{data/eol.json} na każdym pull requeście
\emph{i}~w~miesięcznym odświeżeniu, zanim powstanie jego pull request, więc zmiana po stronie
endoflife.date, która zepsułaby to, na czym polega strona (przemianowany cykl, brakujący
produkt, zniekształcona data), zatrzymuje odświeżenie, zamiast dotrzeć do odwiedzających.
Sprawdza: walidację generatora; że każdy rozpoznawany moniker frameworka (\repopath{v3.5},
\repopath{net40}--\repopath{net481}, \repopath{netcoreapp1.0}--\repopath{netcoreapp3.1},
\repopath{net5.0}--\repopath{net10.0}) trafia w~wymieniony cykl; że dla przykładu .NET
Framework 4.5.1 i~jQuery 1.x są po końcu wsparcia; że formularz ma serwery SQL Server 2012,
2014, 2016, 2017, 2019, 2022 i~Windows Server 2012 R2, 2016, 2019, 2022.

\subsubsection{Testy e2e (Playwright)}

\path{playwright.config.mjs}: katalog \path{e2e/}, lokalny \repopath{baseURL}
\path{http://127.0.0.1:4173/letsgolegacy.legacy-risk-scan/} (albo \repopath{BASE\_URL}),
\repopath{locale} \repopath{pl-PL}, strefa \path{Europe/Warsaw}, ślad zachowywany przy porażce,
zrzut ekranu tylko przy porażce; w~CI \repopath{forbidOnly} i~raport HTML. Dwa projekty:
\repopath{chromium} (zestaw akceptacyjny z~makietą OSV, bez \path{live.spec}, bez ponowień --
zero tolerancji dla niestabilności) i~\repopath{live} (\path{live.spec}, prawdziwe OSV, dwa
ponowienia na wypadek chwilowego błędu sieci). Bez \repopath{BASE\_URL} Playwright uruchamia
\path{scripts/serve.mjs} na porcie 4173 z~bazą \path{/letsgolegacy.legacy-risk-scan/}.

Makieta (\path{e2e/helpers.mjs}) odpowiada za \repopath{api.osv.dev} z~tych samych nagranych
rekordów co testy jednostkowe (NuGet \repopath{Newtonsoft.Json} 6.0.4 →
\repopath{GHSA-5crp-9r3c-p9vr}; npm \repopath{jquery} 1.10.2 → \repopath{GHSA-gxr4-xjj5-5px2},
\repopath{GHSA-rmxg-73gg-4p98}); każde żądanie do innego hosta niż origin strony
i~\repopath{api.osv.dev} jest przerywane i~zapisywane, więc test prywatności na nim zawodzi.
Scenariusze \path{e2e/scan.spec.mjs}:

\begin{itemize}
  \item przykład daje raport z~pozycją po końcu wsparcia i~podatnością z~linkiem;
  \item strona nie pokazuje wymyślonych danych kontaktowych (brak linków \repopath{mailto:}
        i~\repopath{tel:});
  \item prywatność: tylko współrzędne pakietów opuszczają przeglądarkę i~nic nie jest
        zapisywane (sekcja~\ref{sec:risk-scan-prywatnosc});
  \item uszkodzony plik dostaje jasny komunikat po polsku i~żadnego raportu; pusty formularz
        prosi o~plik; plik, który nie jest manifestem, jest nazwany tym, czym jest;
  \item OSV nieosiągalne: wyniki EOL nadal widoczne, z~ostrzeżeniem i~działającym ponowieniem;
  \item serwery wybrane w~formularzu dołączają do raportu bez ponownego pytania OSV;
  \item plik wybrany z~dysku jest czytany lokalnie i~skanowany;
  \item sama klawiatura: Tab dochodzi do „Wypróbuj na przykładzie”, Enter uruchamia skan,
        fokus przechodzi do raportu;
  \item ekran o~szerokości 390~px: bez przewijania w~poziomie, wiersze stają się kartami;
  \item brak poważnych i~krytycznych naruszeń dostępności (axe) przed raportem i~po nim,
        w~trybie jasnym i~ciemnym.
\end{itemize}

\path{e2e/live.spec.mjs} (projekt \repopath{live}) otwiera stronę, uruchamia przykład na
prawdziwym OSV i~sprawdza: brak ostrzeżenia o~OSV, wiersz .NET Framework z~odznaką „Koniec
wsparcia”, link \repopath{GHSA-5crp-9r3c-p9vr} do \path{https://osv.dev/vulnerability/GHSA-5crp-9r3c-p9vr},
że jedynym zewnętrznym originem jest \path{https://api.osv.dev} i~że podatności jest co
najmniej jedna; zapisuje zrzut całej strony (\repopath{LIVE\_SCREENSHOT}).

\subsubsection{Test dymny OSV}

\path{scripts/smoke-osv.mjs} jest jedynym testem, który dowodzi, że jedyna zewnętrzna
zależność strony nadal zachowuje się tak, jak strona zakłada. Sprawdza po kolei: (1) CORS --
preflight \repopath{OPTIONS} dla \repopath{POST} z~originu opublikowanej strony
(\repopath{SMOKE\_ORIGIN}, domyślnie \path{https://konradcinkusz.github.io}) dostaje 2xx
i~nagłówki \path{Access-Control-Allow-*} dla originu, metody \repopath{POST} i~nagłówka
\repopath{Content-Type}, a~sam \repopath{POST} -- \path{Access-Control-Allow-Origin};
(2) kontrakt, przez kod klienta strony -- \repopath{querybatch} znajduje dla NuGet
\repopath{Newtonsoft.Json} 6.0.4 podatność \repopath{GHSA-5crp-9r3c-p9vr}, na której opiera się
przykład, npm \repopath{jquery} 1.10.2 ma co najmniej jedną podatność, a~\path{/v1/vulns}
zwraca rekord ze streszczeniem, wersją z~poprawką i~wagą. Wynik dla
\repopath{newtonsoft.json} pisanego małymi literami jest tylko zapisywany (to na nim opiera się
twierdzenie o~rozróżnianiu wielkości liter). Kody wyjścia: 0 -- sukces albo API zupełnie
nieosiągalne (brak połączenia, limit czasu, 5xx, 429: ostrzeżenie i~\repopath{reachable=false}
w~\repopath{\$GITHUB\_OUTPUT}, żeby dalsze kroki się pominęły); 1 -- OSV odpowiedziało, ale nie
tak, jak strona potrzebuje. Tylko to drugie jest porażką założeń repozytorium.

\zrodla{\path{letsgolegacy.legacy-risk-scan/README.md},
\path{letsgolegacy.legacy-risk-scan/package.json},
\path{letsgolegacy.legacy-risk-scan/test/},
\path{letsgolegacy.legacy-risk-scan/test/eol-data.test.mjs},
\path{letsgolegacy.legacy-risk-scan/playwright.config.mjs},
\path{letsgolegacy.legacy-risk-scan/e2e/helpers.mjs},
\path{letsgolegacy.legacy-risk-scan/e2e/scan.spec.mjs},
\path{letsgolegacy.legacy-risk-scan/e2e/live.spec.mjs},
\path{letsgolegacy.legacy-risk-scan/scripts/smoke-osv.mjs}}

\subsection{Wdrożenie}\label{sec:risk-scan-deploy}

Stronę publikuje GitHub Pages; buduje i~wdraża ją \path{.github/workflows/deploy.yml} przy
każdym pushu do \repopath{main} (i~ręcznie). Grupa współbieżności \repopath{pages} bez
anulowania: w~połowie opublikowana strona jest gorsza niż opóźniona.

\begin{enumerate}
  \item \textbf{\repopath{build}} (limit 20 minut): \texttt{npm ci}, testy jednostkowe,
        instalacja Chromium, \texttt{npm run e2e} (build, potem zestaw e2e z~makietą OSV na
        zbudowanej stronie), \path{actions/configure-pages@v6},
        \path{actions/upload-pages-artifact@v5} z~\path{dist}. Nic nieprzetestowanego nie
        jest publikowane.
  \item \textbf{\repopath{deploy}}: \path{actions/deploy-pages@v5} do środowiska
        \repopath{github-pages}, z~uprawnieniami \texttt{pages: write}
        i~\texttt{id-token: write}; wyjście \repopath{page\_url}.
  \item \textbf{\repopath{verify}} (limit 20 minut): najpierw
        \path{scripts/wait-for-deploy.mjs} czeka, aż strona poda ten commit -- odpytuje
        \path{<page-url>build.json} co 10~sekund (parametr \repopath{t} omija cache, limit
        żądania 15~sekund, domyślny limit całkowity 600~sekund; kod 2 przy złych argumentach,
        1 po przekroczeniu czasu) -- więc kontrola nigdy nie przechodzi ani nie zawodzi na
        poprzednim wdrożeniu. Potem \texttt{npm run e2e:live} z~\repopath{BASE\_URL}
        opublikowanej strony: przykład musi dać raport z~wierszem EOL i~prawdziwą podatnością.
        Zrzut całej strony i~wyniki trafiają do artefaktu \repopath{live-verification}
        (przechowywanego 30 dni). To jest „done when” ticketu L18: publiczna strona zwraca
        raport dla przykładowego pliku.
\end{enumerate}

\textbf{Jednorazowy krok ręczny:} \emph{Settings → Pages → Build and deployment → Source →
GitHub Actions}. Bez niego pierwsze wdrożenie zawodzi na \repopath{configure-pages} z~błędem
404, czego nic w~workflow nie może naprawić.

\zrodla{\path{letsgolegacy.legacy-risk-scan/.github/workflows/deploy.yml},
\path{letsgolegacy.legacy-risk-scan/scripts/wait-for-deploy.mjs},
\path{letsgolegacy.legacy-risk-scan/scripts/build.mjs},
\path{letsgolegacy.legacy-risk-scan/README.md}}

\subsection{CI i~automatyzacja}\label{sec:risk-scan-ci}

\begin{xltabular}{\linewidth}{@{}>{\raggedright\arraybackslash}p{3.1cm} >{\raggedright\arraybackslash}X@{}}
\toprule
Workflow & Wyzwalacze i~joby \\
\midrule
\endfirsthead
\toprule
Workflow & Wyzwalacze i~joby \\
\midrule
\endhead
\path{ci.yml} & Każdy pull request, push do \repopath{main}, co tydzień (cron \texttt{40 6 * * 1} -- OSV może się zmienić pod niezmienionym repozytorium), ręcznie. \repopath{unit}: Node~22, \texttt{npm ci}, \texttt{npm test}. \repopath{e2e}: instalacja Chromium z~zależnościami, \texttt{npm run e2e}; przy porażce artefakt \repopath{e2e-report} (14 dni). \repopath{smoke-osv}: \path{scripts/smoke-osv.mjs}; jeśli OSV jest osiągalne -- \texttt{npm ci}, Chromium, build i~\texttt{npm run e2e:live} na lokalnym buildzie (co dowodzi też CORS z~prawdziwej przeglądarki), artefakt \repopath{smoke-report} (14 dni). Blokujący, z~wyjątkiem całkowitej nieosiągalności OSV \\
\path{deploy.yml} & Push do \repopath{main}, ręcznie; \repopath{build} → \repopath{deploy} → \repopath{verify} (sekcja~\ref{sec:risk-scan-deploy}) \\
\path{update-eol.yml} & Co miesiąc (cron \texttt{17 5 1 * *}), ręcznie oraz na pull requeście zmieniającym \path{scripts/update-eol.mjs}, \path{scripts/lib/eol-source.mjs}, \path{data/eol.json} lub sam workflow. Job \repopath{refresh} (poza pull requestami): regeneruje dane, uruchamia testy jednostkowe na nowych danych i~-- jeśli plik się zmienił -- wypycha gałąź \path{automation/update-eol-data} (commit \texttt{chore(data): refresh end-of-life data from endoflife.date (<data>)}) i~otwiera albo aktualizuje pull request „Odświeżenie danych EOL (endoflife.date, <data>)”, którego opis linkuje przebieg z~testami (pull request otwarty przez \repopath{GITHUB\_TOKEN} nie uruchamia innych workflowów, więc testy biegną tu, zanim PR powstanie). Job \repopath{sync-pr} (na pull requeście): regeneruje z~\path{--keep-date-if-unchanged}, testuje i~commituje wynik na gałąź PR, jeśli treść się różni; na gałęzi z~forka, do której nie da się pisać, zawodzi, gdy dane są nieaktualne. Jedyny zapisujący \path{data/eol.json} \\
\path{codeql.yml} & Pull requesty, push do \repopath{main}, co tydzień (cron \texttt{0 4 * * 1}), ręcznie. Języki \path{javascript-typescript} i~\repopath{actions}, zapytania \path{security-and-quality}, bez kroku build. Uzasadnienie: strona renderuje tekst z~zewnątrz (streszczenia OSV) i~z~pliku odwiedzającego -- dokładnie powierzchnię wstrzyknięć DOM, którą obejmują zapytania JavaScript; workflowy pushują commity i~otwierają pull requesty, co obejmują zapytania dla Actions \\
\path{secret-scan.yml} & Pull requesty, push do \repopath{main}, co tydzień (cron \texttt{30 5 * * 1}), ręcznie; gitleaks na pełnej historii z~obrazu \repopath{:latest} (celowo). Połowa CI dwuczęściowej kontroli; drugą jest hook pre-commit \\
\bottomrule
\end{xltabular}

Job \repopath{refresh} potrzebuje raz ustawienia \emph{Settings → Actions → General →
Workflow permissions → Allow GitHub Actions to create and approve pull requests}; bez niego
gałąź jest wypychana, a~przebieg zawodzi z~komunikatem, jak otworzyć pull request ręcznie.
Commit wypchnięty przez \repopath{GITHUB\_TOKEN} w~\repopath{sync-pr} nie uruchamia ponownie
CI; zrobi to następny push do gałęzi.

Dependabot co tydzień (poniedziałek 06:00) grupuje aktualizacje narzędzi npm (grupa
\repopath{dev-tooling}) i~GitHub Actions. Szablon pull requestu
(\path{.github/pull_request_template.md}) wymaga identyfikatora ticketu z~„done when”
(dosłownie) i~opisu zmiany -- dla tego, co widzi odwiedzający, z~cytatem polskiego tekstu --
oraz dwóch list kontrolnych. Kontrola prywatności (strona obiecuje, że przeglądarkę opuszczają
tylko współrzędne pakietów; zaznacza się jedno):

\begin{checklista}
  \item This change does not touch what the page sends over the network
  \item This change touches it, and the e2e privacy test (\path{e2e/}) still asserts the
        request body
\end{checklista}

Weryfikacja:

\begin{checklista}
  \item \texttt{npm test} (unit)
  \item \texttt{npm run e2e} (Playwright, OSV mocked)
  \item \texttt{./scripts/scan-secrets.sh} clean
  \item If \path{data/eol.json} changed: the diff was read, not only regenerated
\end{checklista}

\zrodla{\path{letsgolegacy.legacy-risk-scan/.github/workflows/ci.yml},
\path{letsgolegacy.legacy-risk-scan/.github/workflows/deploy.yml},
\path{letsgolegacy.legacy-risk-scan/.github/workflows/update-eol.yml},
\path{letsgolegacy.legacy-risk-scan/.github/workflows/codeql.yml},
\path{letsgolegacy.legacy-risk-scan/.github/workflows/secret-scan.yml},
\path{letsgolegacy.legacy-risk-scan/.github/dependabot.yml},
\path{letsgolegacy.legacy-risk-scan/.github/pull_request_template.md}}

\subsection{Znane ograniczenia}\label{sec:risk-scan-ograniczenia}

Według README:

\begin{itemize}
  \item \textbf{OSV dopasowuje identyfikatory NuGet z~rozróżnianiem wielkości liter}
        (sprawdzone na żywym API: test dymny zgłasza \repopath{Newtonsoft.Json} → 1
        podatność, \repopath{newtonsoft.json} → 0). Identyfikatory są wysyłane dokładnie tak,
        jak są zapisane; narzędzia NuGet zapisują kanoniczną wielkość liter, ręcznie
        edytowany \repopath{PackageReference} może jej nie mieć.
  \item Klasyczny \repopath{.csproj} wymienia tylko pakiety, które dostarczają DLL; pakiety
        wyłącznie z~treścią (Bootstrap, większość wtyczek jQuery) widać tylko
        w~\repopath{packages.config} -- raport to mówi.
  \item Ograniczenia (\repopath{composer.json}, \repopath{package.json}, zakresy NuGet) są
        sprawdzane w~najniższej dopuszczalnej wersji; zainstalowana może być nowsza. Pliki lock
        dają dokładne wyniki.
  \item Ręcznie skopiowanych DLL i~komponentów GAC nie da się sprawdzić automatycznie; są
        wymienione jako wymagające ręcznej weryfikacji, a~nie zgadywane.
\end{itemize}

Z~kodu: \repopath{package-lock.json} nie jest obsługiwany (strona prosi o~\repopath{package.json});
właściwość MSBuild zdefiniowana poza wklejonym plikiem nie jest rozwiązywana (np. wersje
z~\path{Directory.Packages.props} trzeba wkleić osobno); raport widzi tylko to, co zapisano
w~jednym pliku.

\zrodla{\path{letsgolegacy.legacy-risk-scan/README.md},
\path{letsgolegacy.legacy-risk-scan/src/lib/manifest.js},
\path{letsgolegacy.legacy-risk-scan/src/lib/parsers/msbuild.js},
\path{letsgolegacy.legacy-risk-scan/src/lib/messages.js}}

\subsection{Status według WORKPLAN}\label{sec:risk-scan-status}

\begin{tabularx}{\linewidth}{@{}>{\raggedright\arraybackslash}p{1.0cm} >{\raggedright\arraybackslash}X >{\raggedright\arraybackslash}X >{\raggedright\arraybackslash}p{4.0cm}@{}}
\toprule
Ticket & Rezultat & Zrobione, gdy & Status \\
\midrule
L18 & Wklejenie \repopath{.csproj} / \repopath{packages.config} / \repopath{composer.json} → raport EOL i~CVE & Publiczna strona zwraca raport dla przykładowego pliku & scalone (\#1, \#2); zrobione, gdy przejdzie job \repopath{verify} w~\path{deploy.yml} -- czeka na krok 1 poniżej \\
\bottomrule
\end{tabularx}

L18 zrealizowano w~dwóch piętrowych pull requestach: \#1 (silnik, pipeline danych EOL,
baseline repozytorium) i~\#2 (strona, wdrożenie, testy e2e). Przed przeglądem sprawdzono
w~CI, że przykład daje w~prawdziwej przeglądarce raport z~wierszami EOL i~prawdziwą
podatnością z~OSV oraz że OSV odpowiada na CORS dla opublikowanego originu. Niesprawdzona
pozostaje sama opublikowana strona, która nie istnieje przed wykonaniem kroków:

\begin{enumerate}
  \item \emph{Settings → Pages → Build and deployment → Source → GitHub Actions} (jednorazowo).
  \item Scalenie \#1, potem \#2 -- wykonane 26~IX~2026. Pierwszy przebieg \path{deploy.yml} na
        \repopath{main} zawiódł na \repopath{configure-pages} z~błędem 404, czego oczekiwano
        przed krokiem 1.
  \item Push do \repopath{main} uruchamia \path{deploy.yml}; jego job \repopath{verify} ładuje
        \textbf{opublikowaną} stronę, uruchamia przykład na prawdziwym OSV i~zachowuje zrzut
        całej strony (artefakt \repopath{live-verification}). Gdy przejdzie, warunek „done when”
        jest spełniony i~L18 można oznaczyć jako zrobiony.
  \item Dla miesięcznego pull requestu z~danymi EOL: \emph{Settings → Actions → General → Allow
        GitHub Actions to create and approve pull requests}.
  \item Gdy ticket L17 dostarczy domenę i~adres kontaktowy, zastąpić placeholder
        \repopath{TODO(L17)} w~\path{src/index.html} (do tego czasu strona nie pokazuje
        wymyślonego adresu).
\end{enumerate}

\zrodla{\path{letsgolegacy.legacy-risk-scan/docs/WORKPLAN.md},
\path{letsgolegacy.legacy-risk-scan/README.md}}
